% Generated draft: requires mathematical and source review.
\documentclass[11pt]{article}
\usepackage[margin=25mm]{geometry}
\usepackage{fontspec}
\setmainfont{Latin Modern Roman}
\usepackage{amsmath,amssymb,mathtools,unicode-math}
\setmathfont{Latin Modern Math}
\usepackage{enumitem,needspace,xurl,hyperref}
\setlength{\parindent}{0pt}
\setlength{\parskip}{5pt}
\setlength{\emergencystretch}{3em}
\newcommand{\N}{\mathbb{N}}
\newcommand{\Z}{\mathbb{Z}}
\newcommand{\Q}{\mathbb{Q}}
\newcommand{\R}{\mathbb{R}}
\newcommand{\C}{\mathbb{C}}
\newcommand{\F}{\mathbb{F}}
\newcommand{\A}{\mathbb{A}}
\newcommand{\PP}{\mathbb{P}}
\newcommand{\E}{\mathbb{E}}
\newcommand{\eps}{\varepsilon}
\newcommand{\dd}{\,\mathrm{d}}
\begin{document}
\section*{UnsolvedMath ID: 30001582}
\textbf{Source problem code:} catalogue entry
\subsection{Definitions and mathematical statement}
A group \(G\) is finitely generated if \(G=\langle X\rangle\) for some finite subset \(X\subseteq G\). It is a torsion group if every element has finite order, that is,
\[
\forall g\in G\ \exists n\in\Z_{\geq 1}\quad g^n=1.
\]
It is residually finite if, for every \(g\in G\setminus\{1\}\), there exist a finite group \(F\) and a homomorphism \(\phi:G\to F\) with \(\phi(g)\neq 1\). For a subgroup \(H\leq G\), write \([G:H]\) for its index.

\textbf{Restricted-subgroup question.} Does there exist an infinite, finitely generated, residually finite torsion group \(G\) such that
\[
\forall H\leq G,\qquad |H|<\infty\ \text{or}\ [G:H]<\infty?
\]
Here the quantifier ranges over all abstract subgroups, including subgroups that are not finitely generated.

\textbf{Arithmetic valuation question.} The source also asks: for a number field \(K\), a prime number \(p\), and a finite set \(S\) of primes, if the group denoted by \(G=G_{K,p,S}\) has positive weighted deficiency, can one choose a valuation \(W\), naturally described by the arithmetic of \(K\), such that
\[
\operatorname{def}_{W}(G_{K,p,S})>0?
\]
The intended quantification is over every such choice of \(K,p,S\) satisfying the positive-weighted-deficiency hypothesis, with an existential choice of \(W\) for that choice of data. A number field means a finite extension of \(\Q\).

The supplied record does not define \(G_{K,p,S}\), specify which primes may belong to \(S\), or define the admissible valuations and the invariant \(\operatorname{def}_{W}\). It also gives no mathematical criterion for a valuation to be ``naturally described by the arithmetic of \(K\).'' Consequently the arithmetic question is retained as a qualified source-level question; its hypotheses and requested conclusion cannot be fully formalized from this record alone.
\subsection{Short English statement}
The first question asks for an infinite finitely generated residually finite torsion group whose every subgroup is finite or has finite index. The second asks whether positive weighted deficiency for the arithmetic group \(G_{K,p,S}\) can be witnessed by an arithmetically natural valuation. The source does not supply enough definitions to make the second question fully precise. Its supplied partial-solution assessment concerns finitely generated subgroups and does not replace the first question's quantifier over all subgroups.
\subsection{Sources}
\begin{itemize}
\item Local UnsolvedMath registry: supplied complete source record, permanent ID \(30001582\), title ``Residually Finite Torsion Groups with Restricted Subgroups,'' source problem code \texttt{catalogue entry}, in the Oberwolfach Reports Open Problems collection. The record lists proposed year \(2010\) and supplies no individual proposer.
\item Supplied authoritative source citation: \emph{Actions and Invariants of Residually Finite Groups: Asymptotic Methods} (2011), Oberwolfach Reports, pp.\ 2335--2391, DOI \texttt{10.4171/owr/2010/40}. Original OWR report: \url{https://doi.org/10.4171/owr/2010/40}. The supplied background describes these reports as documenting mathematical workshops and their open-problem discussions.
\item Supplied literature reference: Ershov, \emph{Golod--Shafarevich groups survey}, \url{https://m-ershov.github.io/Research/gssurvey_revised.pdf}.
\item Literature assessment dated exactly \texttt{2026-08-21}: the record labels the problem ``partially solved'' and states that the all-subgroups residually finite torsion-group question remains open, that the positive weighted-deficiency construction proves the analogous property for every finitely generated subgroup, and that the valuation question remains open. This is supplied dated curation derived from the OWR curation, not an independently verified literature assessment; time-sensitive status and cited claims require independent verification.
\item The registry labels the statement \texttt{corrected\_verified} and reports that its mathematical content was checked against the official Oberwolfach report, with notation, context, and question form normalized without changing the source-level claim. This is a supplied verification assertion, not verification performed for this draft.
\end{itemize}
\end{document}
